% ============================================================
%  Презентация для защиты магистерской диссертации
%  НИУ ВШЭ — Факультет компьютерных наук
%  Автор: Семенюченко А. Г. | Науч. рук.: Паточенко Е. А.
%  Тема: Прогнозирование ВРП регионов РФ (ансамблевые ML)
%
%  Компиляция: xelatex presentation_hse.tex  (дважды)
% ============================================================

\documentclass[aspectratio=169,10pt]{beamer}

% ─── XeLaTeX + Russian ───────────────────────────────────────────────────────
\usepackage{fontspec}
\usepackage{polyglossia}
\setmainlanguage{russian}
\setotherlanguage{english}

\setmainfont{Helvetica Neue}
\setsansfont{Helvetica Neue}
\setmonofont{Courier New}
\newfontfamily\cyrillicfontsf{Helvetica Neue}
\newfontfamily\cyrillicfonttt{Courier New}

% ─── Пакеты ──────────────────────────────────────────────────────────────────
\usepackage{graphicx}
\usepackage{booktabs}
\usepackage{xcolor}
\usepackage{tikz}
\usetikzlibrary{positioning, shapes.geometric, calc}
\usepackage{amsmath, amssymb}
\usepackage{array}
\usepackage{colortbl}
\usepackage{tabularx}
\usepackage{microtype}

% ─── Цвета брендбука ВШЭ ────────────────────────────────────────────────────
\definecolor{hsedark}{RGB}{15,45,105}       % Тёмно-синий (Pantone 287)
\definecolor{hseblue}{RGB}{35,75,155}       % Синий (Pantone 2728)
\definecolor{hsered}{RGB}{230,30,60}        % Красный акцент (HSE brand)
\definecolor{hsegray}{RGB}{230,230,230}
\definecolor{hsemidgray}{RGB}{191,191,191}
\definecolor{hsedarkgray}{RGB}{128,128,128}
\definecolor{hsegreen}{RGB}{0,155,100}
\definecolor{hseorange}{RGB}{235,105,30}
\definecolor{hsetablehl}{RGB}{205,220,240}

% ─── Beamer: цвета ───────────────────────────────────────────────────────────
\setbeamercolor{structure}{fg=hsedark}
\setbeamercolor{background canvas}{bg=white}
\setbeamercolor{frametitle}{fg=white, bg=hsedark}
\setbeamercolor{title}{fg=hsedark}
\setbeamercolor{subtitle}{fg=hseblue}
\setbeamercolor{author}{fg=hsedark}
\setbeamercolor{date}{fg=hsedark}
\setbeamercolor{normal text}{fg=black}
\setbeamercolor{block title}{fg=white, bg=hsedark}
\setbeamercolor{block body}{fg=black, bg=hsegray}
\setbeamercolor{block title alerted}{fg=white, bg=hsered}
\setbeamercolor{block body alerted}{fg=black, bg=hsered!8}
\setbeamercolor{alerted text}{fg=hsered}
\setbeamercolor{example text}{fg=hsegreen!70!black}
\setbeamercolor{itemize item}{fg=hsedark}
\setbeamercolor{itemize subitem}{fg=hseblue}
\setbeamercolor{enumerate item}{fg=hsedark}
\setbeamercolor{footline}{bg=hsedark}

% ─── Beamer: шрифты ──────────────────────────────────────────────────────────
\setbeamerfont{frametitle}{size=\large, series=\bfseries}
\setbeamerfont{title}{size=\Large, series=\bfseries}
\setbeamerfont{subtitle}{size=\normalsize}
\setbeamerfont{author}{size=\normalsize}
\setbeamerfont{date}{size=\small}
\setbeamerfont{footline}{size=\tiny}
\setbeamerfont{block title}{size=\normalsize, series=\bfseries}
\setbeamerfont{section in toc}{series=\bfseries}

% ─── Убираем навигацию ───────────────────────────────────────────────────────
\setbeamertemplate{navigation symbols}{}

% ─── Маркеры списков ─────────────────────────────────────────────────────────
\setbeamertemplate{itemize item}{%
  \tikz\fill[hsedark] (0,0) -- (0.2em,0.35em) -- (-0.2em,0.35em) -- cycle;}
\setbeamertemplate{itemize subitem}{%
  \tikz\fill[hseblue] (0,0) rectangle (0.35em,0.35em);}
\setbeamertemplate{itemize subsubitem}{%
  \tikz\fill[hsedarkgray] circle (0.18em);}
\setbeamertemplate{enumerate item}{%
  \textbf{\color{hsedark}\insertenumlabel.}}

% ─── Заголовок слайда ────────────────────────────────────────────────────────
\setbeamertemplate{frametitle}{%
  \vspace{-0.02cm}%
  \begin{beamercolorbox}[%
      wd=\paperwidth, ht=0.75cm, dp=0.20cm,
      leftskip=0.60cm, rightskip=0.40cm]{frametitle}%
    \usebeamerfont{frametitle}%
    \usebeamercolor[fg]{frametitle}%
    \insertframetitle%
  \end{beamercolorbox}%
}

% ─── Нижний колонтитул ───────────────────────────────────────────────────────
\setbeamertemplate{footline}{%
  \begin{beamercolorbox}[%
      wd=\paperwidth, ht=0.48cm, dp=0.16cm,
      leftskip=0.40cm, rightskip=0.40cm]{footline}%
    {\color{white}\tiny
      НИУ «Высшая школа экономики» \textbar{} ФКН \textbar{} ОП «Анализ больших данных»}%
    \hfill
    {\color{white}\tiny
      Семенюченко~А.\,Г. \quad\textbar\quad Прогнозирование ВРП субъектов РФ}%
    \hfill
    {\color{white}\tiny \insertframenumber\,/\,\inserttotalframenumber\hspace{0.3cm}}%
  \end{beamercolorbox}%
}

% ─── Верхняя полоска ─────────────────────────────────────────────────────────
\makeatletter
\setbeamertemplate{headline}{%
  \ifbeamer@plainframe\else%
  \begin{beamercolorbox}[wd=\paperwidth,ht=0.04cm,dp=0pt]{footline}%
  \end{beamercolorbox}%
  \fi%
}
\makeatother

% ─── Путь к изображениям ─────────────────────────────────────────────────────
\graphicspath{{/Users/tigran/python_projects_epta/personal_projects/disser_epta/grp_forecast/figures/}}

% ─── Метаданные ──────────────────────────────────────────────────────────────
\title{Прогнозирование ВРП субъектов Российской Федерации}
\subtitle{с помощью ансамблевых методов машинного обучения}
\author{Семенюченко~А.\,Г.}
\institute{НИУ ВШЭ — Факультет компьютерных наук}
\date{Москва, 2026}

% ─── TOC-стиль ───────────────────────────────────────────────────────────────
\setbeamertemplate{section in toc}{%
  \textcolor{hsedark}{\textbf{\inserttocsectionnumber.}}\ \inserttocsection}

%==============================================================================
\begin{document}
%==============================================================================

% ══════════════════════════════════════════════════════════════════════════════
% СЛАЙД 1: Титульный
% ══════════════════════════════════════════════════════════════════════════════
\begin{frame}[plain]
\begin{tikzpicture}[remember picture, overlay]
  \fill[hsedark]
    (current page.north west) rectangle
    ([yshift=-2.8cm]current page.north east);
  % Красная акцентная полоса (брендбук ВШЭ)
  \fill[hsered]
    ([yshift=-2.80cm]current page.north west) rectangle
    ([yshift=-2.98cm]current page.north east);
  \fill[hsedark]
    (current page.south west) rectangle
    ([yshift=0.72cm]current page.south east);
  \draw[white, line width=1.8pt]
    ([xshift=1.20cm, yshift=-1.40cm]current page.north west) circle (0.62cm);
  \node[white, font=\bfseries\footnotesize] at
    ([xshift=1.20cm, yshift=-1.40cm]current page.north west) {ВШЭ};
  \node[white, font=\scriptsize, anchor=west] at
    ([xshift=2.10cm, yshift=-0.98cm]current page.north west)
    {НАЦИОНАЛЬНЫЙ ИССЛЕДОВАТЕЛЬСКИЙ УНИВЕРСИТЕТ};
  \node[white, font=\scriptsize\bfseries, anchor=west] at
    ([xshift=2.10cm, yshift=-1.50cm]current page.north west)
    {«ВЫСШАЯ ШКОЛА ЭКОНОМИКИ»};
  \node[white, font=\tiny, anchor=west] at
    ([xshift=2.10cm, yshift=-2.10cm]current page.north west)
    {Факультет компьютерных наук \textbar{}
     Образовательная программа «Анализ больших данных»};
  \node[white, font=\scriptsize\bfseries, anchor=west] at
    ([xshift=0.50cm, yshift=0.46cm]current page.south west)
    {Магистерская диссертация};
  \node[white, font=\scriptsize, anchor=east] at
    ([xshift=-0.50cm, yshift=0.46cm]current page.south east)
    {Москва, 2026};
\end{tikzpicture}

\vspace{2.6cm}
\begin{center}
  {\usebeamerfont{title}\color{hsedark}%
    Прогнозирование валового регионального\\[0.2em]%
    продукта (ВРП) субъектов Российской Федерации}\\[0.4em]
  {\color{hseblue}\large
    с помощью ансамблевых методов машинного обучения}\\[0.5em]
  {\color{hsemidgray}\rule{0.65\textwidth}{0.4pt}}\\[0.4em]
  {\small
  \begin{tabular}{r@{\hspace{0.5em}}l}
    \textbf{\color{hsedark}Выполнил:}           &
      Семенюченко Артём Геннадьевич, студент ФКН НИУ ВШЭ\\[0.2em]
    \textbf{\color{hsedark}Науч. руководитель:} &
      Паточенко Евгений Анатольевич, ст. преп. ФКН НИУ ВШЭ
  \end{tabular}}
\end{center}
\end{frame}

% ══════════════════════════════════════════════════════════════════════════════
% СЛАЙД 2: Содержание
% ══════════════════════════════════════════════════════════════════════════════
\begin{frame}{Структура доклада}
  \tableofcontents[subsectionstyle=hide]
\end{frame}

%==============================================================================
\section{Актуальность и цель исследования}
%==============================================================================

% ══════════════════════════════════════════════════════════════════════════════
% СЛАЙД 3: Актуальность
% ══════════════════════════════════════════════════════════════════════════════
\begin{frame}{Актуальность исследования}
\begin{columns}[T]
  \begin{column}{0.55\textwidth}
    \begin{block}{Почему прогнозирование ВРП регионов?}
      \begin{itemize}\setlength{\itemsep}{0pt}
        \item ВРП~— ключевой индикатор развития 85~субъектов~РФ
        \item \textbf{Территориальная неоднородность}: сырьевые,
              индустриальные, аграрные, бюджетно-зависимые регионы
        \item \textbf{Лаг публикации} Росстата~— бюджетное планирование
              требует оперативных оценок до выхода официальных данных
      \end{itemize}
    \end{block}
  \end{column}
  \begin{column}{0.43\textwidth}
    \begin{block}{Вызовы задачи}
      \begin{itemize}\setlength{\itemsep}{0pt}
        \item Панельная структура: (регион, год)
        \item Временна\'{я} зависимость + межрегиональная неоднородность
        \item Риск утечки данных (\textit{leakage})
        \item Структурные разрывы: 2008–09, 2014–15, 2020, 2022–23
        \item Малая панель ($\approx$\,20 набл./регион)
      \end{itemize}
    \end{block}
  \end{column}
\end{columns}
\vspace{0.1em}
\begin{alertblock}{}
  \centering
  \textbf{Цель}: воспроизводимый pipeline прогнозирования ВРП на основе ML
  с корректной out-of-sample валидацией и явным контролем leakage
\end{alertblock}
\end{frame}

% ══════════════════════════════════════════════════════════════════════════════
% СЛАЙД 4: Цель и задачи
% ══════════════════════════════════════════════════════════════════════════════
\begin{frame}{Цель и задачи исследования}
\begin{block}{Цель работы}
  Разработка и оценка воспроизводимого пайплайна прогнозирования ВРП
  субъектов~РФ на панельных временны\'{х} данных с использованием
  эконометрических базовых моделей и ансамблевых методов машинного обучения.
\end{block}
\vspace{0.25em}
\textbf{\color{hsedark}Ключевые задачи:}
\begin{enumerate}
  \item Сформировать панельную базу данных за \textbf{2005–2024~гг.}
        (Росстат, ЕМИСС, Банк России, Минэкономразвития, Казначейство)
  \item Построить \textbf{базовые} (наивный, AR(1), Ridge с FE) и
        \textbf{ансамблевые} модели (LightGBM, CatBoost, GPBoost, EmbMLP)
  \item Реализовать leakage-free \textbf{expanding-window} кросс-валидацию
        (тестовый период: 2015–2023)
  \item Провести \textbf{ablation study} по блокам признаков и
        \textbf{robustness checks} (выбросы и кризисные периоды)
  \item Построить \textbf{SHAP-интерпретацию} лучшей одиночной модели
        и сформировать взвешенный ансамбль
\end{enumerate}
\end{frame}

%==============================================================================
\section{Данные и методология}
%==============================================================================

% ══════════════════════════════════════════════════════════════════════════════
% СЛАЙД 5: Данные и признаковое пространство
% ══════════════════════════════════════════════════════════════════════════════
\begin{frame}{Данные и признаковое пространство}
\begin{columns}[T]
  \begin{column}{0.47\textwidth}
    \textbf{\color{hsedark}Эмпирическая база:}
    \begin{itemize}
      \item 85~субъектов РФ, \textbf{2005–2024~гг.}
      \item Целевая переменная:
            $\Delta\log\text{ВРП}_{i,t} = \log\text{ВРП}_{i,t} - \log\text{ВРП}_{i,t-1}$
      \item Горизонт прогноза: $h=1$ (один год вперёд)
    \end{itemize}
    \vspace{0.3em}
    \textbf{\color{hsedark}Предобработка:}
    \begin{itemize}
      \item Дефлирование ИПЦ (денежные показатели)
      \item \texttt{RobustScaler} — квартильная нормировка,
            устойчивость к выбросам (Москва, ХМАО, ЯНАО)
      \item Все операции — \textbf{только на Train}$_k$ каждого fold'а
    \end{itemize}
  \end{column}
  \begin{column}{0.52\textwidth}
    \textbf{\color{hsedark}Блоки признаков:}\\[0.2em]
    {\scriptsize$X_{i,t} {=} \bigl[X^{\text{lag}},X^{\text{prod}},X^{\text{inv}},
     X^{\text{lab}},X^{\text{bud}},X^{\text{macro}},X^{\text{price}},
     X^{\text{shock}},X^{\text{cat}}\bigr]$}
    \vspace{0.1em}
    {\scriptsize
    \begin{tabular}{@{}lp{3.6cm}@{}}
      \toprule
      Блок & Содержание \\
      \midrule
      \textsc{lag}   & Лаги ВРП, скользящие средние \\
      \textsc{prod}  & Пром-сть, с/х, стр-во, торговля \\
      \textsc{inv}   & Инвестиции в осн. капитал \\
      \textsc{lab}   & Занятость, зарплата, доходы \\
      \textsc{bud}   & Доходы/расходы бюджета \\
      \textsc{macro} & Ключевая ставка, кредиты \\
      \textsc{price} & USD/RUB волат., прод. корзина \\
      \textsc{shock} & Dummies кризисов (2009, 2015, 2020) \\
      \textsc{cat}   & Тип региона, фед. округ \\
      \bottomrule
    \end{tabular}}
  \end{column}
\end{columns}
\end{frame}

% ══════════════════════════════════════════════════════════════════════════════
% СЛАЙД 6: Схема валидации
% ══════════════════════════════════════════════════════════════════════════════
\begin{frame}{Методология: expanding-window валидация и метрики}
\begin{columns}[T]
  \begin{column}{0.54\textwidth}
    \textbf{\color{hsedark}Схема расширяющегося окна (9 fold'ов):}
    \vspace{0.2em}
    \begin{center}
    \begin{tikzpicture}[xscale=0.42, yscale=0.65]
      \draw[->, thick, color=hsedark] (-0.1,0) -- (16.8,0);
      \foreach \x/\yr in {0/2005, 5.45/2010, 9.9/2015, 14.4/2020, 15.8/2023}
        \node[font=\tiny, below=1pt] at (\x,0) {\yr};

      \fill[hseblue!35] (0,0.25) rectangle (9.90,0.75);
      \fill[hsegreen!70] (9.90,0.25) rectangle (10.80,0.75);
      \node[font=\tiny, left=2pt] at (0,0.50) {fold 1};

      \fill[hseblue!35] (0,0.90) rectangle (10.80,1.40);
      \fill[hsegreen!70] (10.80,0.90) rectangle (11.70,1.40);
      \node[font=\tiny, left=2pt] at (0,1.15) {fold 2};

      \fill[hseblue!35] (0,1.55) rectangle (11.70,2.05);
      \fill[hsegreen!70] (11.70,1.55) rectangle (12.60,2.05);
      \node[font=\tiny, left=2pt] at (0,1.80) {fold 3};

      \node at (8, 2.40) {\Large\ldots};

      \fill[hseblue!35] (0,2.80) rectangle (14.40,3.30);
      \fill[hsegreen!70] (14.40,2.80) rectangle (15.30,3.30);
      \node[font=\tiny, left=2pt] at (0,3.05) {fold 9};

      % Легенда — внутри диаграммы, над fold 9
      \fill[hseblue!35] (3.0,3.65) rectangle (4.5,3.90);
      \node[right, font=\tiny] at (4.5,3.78) {Train};
      \fill[hsegreen!70] (8.5,3.65) rectangle (10.0,3.90);
      \node[right, font=\tiny] at (10.0,3.78) {Test (1 год)};
    \end{tikzpicture}
    \end{center}
    \vspace{0.1em}
    {\small\color{hsedark}Начальное окно: 2005–2014 ($\geq\!10$~лет).
    Гиперпараметры~— Optuna (60~trials, только на Train$_k$).}
  \end{column}
  \begin{column}{0.44\textwidth}
    \textbf{\color{hsedark}Постановка задачи:}
    \[\hat{y}_{i,t+1} = f_\theta\!\left(X_{i,t}^{\text{available}}\right)\]
    \[y_{i,t} = \Delta\log\bigl(\text{GRP}_{i,t}^{\text{real}}\bigr)\]

    \vspace{0.3em}
    \textbf{\color{hsedark}Метрики качества:}
    \begin{align*}
      &\text{RMSE} = \sqrt{\tfrac{1}{n}\textstyle\sum_j(y_j-\hat{y}_j)^2}\\[3pt]
      &\text{MAE}  = \tfrac{1}{n}\textstyle\sum_j|y_j-\hat{y}_j|\\[3pt]
      &R^2_{\text{OOS}} = 1 -
        \frac{\sum_j(y_j-\hat{y}_j)^2}{\sum_j(y_j-\hat{y}_j^{\text{bench}})^2}
    \end{align*}
    {\small $R^2_{\text{OOS}}>0$ — превосходство над наивным прогнозом.}
  \end{column}
\end{columns}
\end{frame}

% ══════════════════════════════════════════════════════════════════════════════
% СЛАЙД 7: Модели
% ══════════════════════════════════════════════════════════════════════════════
\begin{frame}{Сравниваемые модели}
\begin{columns}[T]
  \begin{column}{0.48\textwidth}
    \begin{block}{Базовые модели}
      \begin{itemize}
        \item \textbf{Наивный}: $\hat{y}_{i,t+1} = y_{i,t}$
        \item \textbf{AR(1)}: $y_{i,t+1} = \alpha_i + \rho_i\,y_{i,t} + \varepsilon$
              (по каждому региону)
        \item \textbf{Ridge с~FE}: L2-регрессия с индивидуальными
              и временны\'{м}и фиксированными эффектами
      \end{itemize}
    \end{block}
    \vspace{0.2em}
    \begin{block}{Взвешенный ансамбль}
      {\small
      $\hat{y}^{\text{ens}} = w\,\hat{y}^{\text{CB}} + (1-w)\,\hat{y}^{\text{LGB}}$\\[3pt]
      Перебор $w\in[0{,}00;\,1{,}00]$, шаг $0{,}05$, min pooled RMSE\\[3pt]
      \textbf{Оптимум: $w^*=0{,}25$}~$\Rightarrow$~CB$\!\times\!$0,25 + LGB$\!\times\!$0,75}
    \end{block}
  \end{column}
  \begin{column}{0.50\textwidth}
    \begin{block}{Ансамблевые алгоритмы ML}
      \begin{itemize}
        \item \textbf{LightGBM}: градиентный бустинг, leaf-wise рост, EFB
        \item \textbf{CatBoost}: ordered boosting, симметричные деревья,
              устойчивость к кат. признакам
        \item \textbf{GPBoost}: бустинг + случайные эффекты регионов\\[-2pt]
              {\small$y_{i,t} {=} F(X_{i,t}){+}b_i{+}\varepsilon_{i,t},\;
               b_i{\sim}\mathcal{N}(0,\sigma_b^2)$}
        \item \textbf{EmbMLP}: нейросеть с обучаемыми эмбеддингами регионов
      \end{itemize}
    \end{block}
  \end{column}
\end{columns}
\end{frame}

%==============================================================================
\section{Результаты экспериментов}
%==============================================================================

% ══════════════════════════════════════════════════════════════════════════════
% СЛАЙД 8: Поиск веса ансамбля
% ══════════════════════════════════════════════════════════════════════════════
\begin{frame}{Взвешенный ансамбль: поиск оптимального веса}
\begin{columns}[T]
  \begin{column}{0.56\textwidth}
    \includegraphics[width=\linewidth]{ensemble_weight_search.png}
  \end{column}
  \begin{column}{0.42\textwidth}
    \vspace{0.4em}
    \textbf{\color{hsedark}Процедура post-hoc оптимизации:}
    \begin{itemize}
      \item После независимой CV каждой подмодели
      \item Перебор $w \in \{0{,}00;\,0{,}05;\,\ldots;\,1{,}00\}$
      \item Минимизация pooled RMSE на hold-out 9~fold'ов
      \item Leakage исключён — Training target не задействуется
    \end{itemize}
    \vspace{0.3em}
    \begin{block}{Результат}
      $w^* = 0{,}25$;\quad $\text{RMSE}(w^*) = 0{,}0382$\\[2pt]
      LightGBM вносит бо\'ольший вклад (доля 0,75).\\[2pt]
      Выигрыш ансамбля над LGB~--- \textbf{1,6\%} RMSE.
    \end{block}
  \end{column}
\end{columns}
\end{frame}

% ══════════════════════════════════════════════════════════════════════════════
% СЛАЙД 9: Сводная таблица метрик
% ══════════════════════════════════════════════════════════════════════════════
\begin{frame}{Сравнение моделей: сводная таблица метрик}
\begin{columns}[T]
  \begin{column}{0.56\textwidth}
    \includegraphics[width=\linewidth]{metrics_table.png}
  \end{column}
  \begin{column}{0.42\textwidth}
    \vspace{0.2em}
    \begin{block}{\small Ключевые итоги}
      {\small
      \begin{itemize}\setlength{\itemsep}{0pt}
        \item \textbf{Лучшая}: ансамбль CB$\!\times\!$0,25+LGB$\!\times\!$0,75\\
              RMSE~$=\mathbf{0{,}0382}$, снижение \textbf{8,8\%} vs наивный;
              $R^2_{\text{OOS}}=0{,}144>0$
        \item \textbf{LightGBM}: RMSE~$=0{,}0388$, лучшая одиночная
        \item \textbf{CatBoost}: RMSE~$=0{,}0408$
        \item AR(1), EmbMLP, Ridge~--- \textit{уступают наивному}
        \item \textbf{GPBoost}: нестабилен (OOS отключает GP)
      \end{itemize}}
    \end{block}
    \vspace{0.1em}
    {\footnotesize\color{hsedark}MAPE~$=\mathbf{2{,}88\%}$;\;
      MAE~$=6{,}24$~млрд руб. (vs~7,25 наивный)}
  \end{column}
\end{columns}
\end{frame}

% ══════════════════════════════════════════════════════════════════════════════
% СЛАЙД 10: RMSE по времени
% ══════════════════════════════════════════════════════════════════════════════
\begin{frame}{RMSE прогнозов: тестовые годы 2015–2023}
\begin{center}
  \includegraphics[width=0.85\linewidth]{rmse_and_growth_combined.png}
\end{center}
\vspace{-0.7em}
\begin{columns}
  \begin{column}{0.50\textwidth}
    {\small\textbf{\color{hsedark}Обучающий период (2006–2014):}\\
     Ср. рост ВРП: 4,2\%/год;\;
     макс.~9,7\% (2007);\; мин.~$-$4,5\% (2009)}
  \end{column}
  \begin{column}{0.48\textwidth}
    {\small\textbf{\color{hsedark}Тестовый период (2015–2023):}\\
     Ансамбль CB+LGB сохраняет \textbf{наименьший RMSE}
     во все кризисные годы (2020, 2022–23)}
  \end{column}
\end{columns}
\end{frame}

% ══════════════════════════════════════════════════════════════════════════════
% СЛАЙД 11: Ablation study
% ══════════════════════════════════════════════════════════════════════════════
\begin{frame}{Ablation study: вклад блоков признаков (LightGBM)}
\begin{columns}[T]
  \begin{column}{0.56\textwidth}
    \includegraphics[width=\linewidth]{ablation_study.png}
  \end{column}
  \begin{column}{0.42\textwidth}
    \vspace{0.3em}
    \textbf{\color{hsedark}Сценарии последовательного добавления:}
    \vspace{0.2em}
    {\small
    \begin{tabular}{@{}lp{3.6cm}@{}}
      \toprule
      Сценарий & Состав \\
      \midrule
      \textbf{A0} & Только лаги ВРП \\
      \textbf{A1} & A0 + производственный блок \\
      \textbf{A2} & A1 + инвестиц. + трудовой + бюджетный \\
      \textbf{A3} & A2 + макро + цены + кат. + кризисные \\
      \bottomrule
    \end{tabular}}
    \vspace{0.3em}
    \begin{block}{Вывод}
      \textbf{Наибольший прирост}: A0$\to$A1 (производственный блок),
      RMSE: $0{,}044\to0{,}043$. A2 и A3 $\approx$ A1.
    \end{block}
  \end{column}
\end{columns}
\end{frame}

% ══════════════════════════════════════════════════════════════════════════════
% СЛАЙД 12: Robustness checks
% ══════════════════════════════════════════════════════════════════════════════
\begin{frame}{Анализ устойчивости (Robustness Checks)}
\begin{columns}[T]
  \begin{column}{0.50\textwidth}
    \textbf{\color{hsedark}RC1~— без регионов-выбросов}\\[1pt]
    {\small Исключены: Москва, ХМАО, ЯНАО, Сахалинская обл.}
    \vspace{0.1em}
    \includegraphics[width=\linewidth]{rc1_outliers.png}
    \vspace{0.05em}
    {\small\color{hsedark}$\Rightarrow$ Ансамбль сохраняет лидерство;
     LightGBM опережает CatBoost~— ранжирование устойчиво}
  \end{column}
  \begin{column}{0.48\textwidth}
    \textbf{\color{hsedark}RC2~— кризисные vs нормальные годы}\\[1pt]
    {\small Кризисные fold'ы: 2008–09, 2014–15, 2020, 2022–23}
    \vspace{0.1em}
    \includegraphics[width=\linewidth]{rc2_crisis.png}
    \vspace{0.05em}
    {\small\color{hsedark}$\Rightarrow$ Ансамбль лучший в обеих подгруппах.
     Бустинг устойчив в кризисные периоды (Lashina \& Grishunin, 2023)}
  \end{column}
\end{columns}
\end{frame}

%==============================================================================
\section{SHAP-интерпретация}
%==============================================================================

% ══════════════════════════════════════════════════════════════════════════════
% СЛАЙД 13: SHAP global
% ══════════════════════════════════════════════════════════════════════════════
\begin{frame}{SHAP-интерпретация: глобальный анализ (LightGBM)}
\begin{columns}[T]
  \begin{column}{0.56\textwidth}
    \includegraphics[width=\linewidth]{shap_global.png}
  \end{column}
  \begin{column}{0.42\textwidth}
    \vspace{0.2em}
    \textbf{\color{hsedark}Топ-5 признаков по $\overline{|\phi_j|}$:}
    \begin{enumerate}
      \small
      \item {\footnotesize\ttfamily usd\_rub\_volatility\_lag1}\\[-1pt]
            {\tiny волатильность USD/RUB; вклад в~2~раза выше след.}
      \item {\footnotesize\ttfamily min\_food\_basket\_cost\_lag1}\\[-1pt]
            {\tiny стоимость минимального набора продуктов}
      \item {\footnotesize\ttfamily policy\_rate\_lag1}\\[-1pt]
            {\tiny номинальная ключевая ставка ЦБ РФ}
      \item {\footnotesize\ttfamily policy\_rate\_real\_lag1}\\[-1pt]
            {\tiny реальная ключевая ставка}
      \item {\footnotesize\ttfamily paid\_services\_to\_grp\_lag1}\\[-1pt]
            {\tiny доля платных услуг в ВРП}
    \end{enumerate}
    \vspace{0.2em}
    \begin{block}{Ключевой вывод}
      {\small\textbf{Монетарные переменные} доминируют; USD/RUB — главный фактор риска.}
    \end{block}
  \end{column}
\end{columns}
\end{frame}

% ══════════════════════════════════════════════════════════════════════════════
% СЛАЙД 14: SHAP по типам регионов
% ══════════════════════════════════════════════════════════════════════════════
\begin{frame}{SHAP: вклад блоков признаков по типам регионов}
\begin{columns}[T]
  \begin{column}{0.56\textwidth}
    \includegraphics[width=\linewidth]{shap_group_heatmap.png}
  \end{column}
  \begin{column}{0.42\textwidth}
    \vspace{0.3em}
    \textbf{\color{hsedark}Интерпретация по блокам:}
    \vspace{0.2em}
    {\footnotesize
    \begin{tabular}{@{}lll@{}}
      \toprule
      Блок & Вклад & Вывод \\
      \midrule
      \textsc{macro}  & 0,21–0,27 & Лидер везде \\
      \textsc{price}  & 0,16–0,21 & 2-й везде \\
      \textsc{prod}   & 0,17 (сырьев.) & Важен для добычи \\
      \textsc{inv}    & 0,11–0,16 & Выше у сервисных \\
      \textsc{cat}    & $\approx 0$ & Нулевой вклад \\
      \textsc{lag}    & 0,09–0,10 & Низкая автокорр. \\
      \bottomrule
    \end{tabular}}
    \vspace{0.2em}
    \begin{alertblock}{\small Вывод}
      {\small Все типы регионов однородно реагируют на монетарные шоки ДКП ЦБ РФ.}
    \end{alertblock}
  \end{column}
\end{columns}
\end{frame}

%==============================================================================
\section{Заключение}
%==============================================================================

% ══════════════════════════════════════════════════════════════════════════════
% СЛАЙД 15: Анализ ошибок + тест DM
% ══════════════════════════════════════════════════════════════════════════════
\begin{frame}{Анализ ошибок и тест Дибольда–Мариано}
\begin{columns}[T]
  \begin{column}{0.50\textwidth}
    \textbf{\color{hsedark}Тепловая карта ошибок (MAE, регион × год)}
    \vspace{0.1em}
    \includegraphics[width=\linewidth]{error_heatmap.png}
    {\tiny Наибольшие ошибки~--- регионы с нестационарной структурой
     эк-ки (Ненецкий АО, Калининградская обл., Калмыкия)}
  \end{column}
  \begin{column}{0.48\textwidth}
    \textbf{\color{hsedark}Тест Дибольда–Мариано (DM)}\\[1pt]
    $H_0$: обе модели одинаково точны\\
    $d_t = e_{1,t}^2 - e_{2,t}^2$;\quad $\text{DM} = \bar{d}/\text{SE}(\bar{d})$
    \vspace{0.3em}
    {\scriptsize
    \begin{tabular}{@{}lrcc@{}}
      \toprule
      Сравнение & DM & $p_{\text{дв}}$ & $p_{\text{одн}}$ \\
      \midrule
      CatBoost vs Наивный & 0,046 & 0,963 & 0,482 \\
      \rowcolor{hsetablehl}
      \textbf{LightGBM vs Наивный} & \textbf{6,247} &
        $<\!0{,}001$ & $<\!0{,}001$ \\
      \bottomrule
    \end{tabular}}
    \vspace{0.2em}
    \begin{block}{\small Вывод}
      {\small LightGBM \textbf{статистически достоверно} превосходит
      наивный прогноз ($\alpha=0{,}001$, DM$=6{,}247$).\\[1pt]
      CatBoost~--- выигрыш в RMSE есть, но не значим.\\[1pt]
      Ансамбль наследует значимость от LGB (доля 0,75).}
    \end{block}
  \end{column}
\end{columns}
\end{frame}

% ══════════════════════════════════════════════════════════════════════════════
% СЛАЙД 16: Заключение
% ══════════════════════════════════════════════════════════════════════════════
\begin{frame}{Заключение: основные результаты}
\begin{columns}[T]
  \begin{column}{0.55\textwidth}
    \begin{block}{Главный результат}
      \textbf{Ансамбль CB$\!\times\!$0,25 + LGB$\!\times\!$0,75} превзошёл все 8 моделей:
      \begin{itemize}
        \item RMSE~$= \mathbf{0{,}0382}$~---
              снижение \textbf{8,8\%} vs наивный ($0{,}0419$)
        \item $R^2_{\text{OOS}} = 0{,}144 > 0$~---
              устойчивое превосходство OOS
        \item \textbf{MAPE~$= 2{,}88\%$}, MAE~$= 6{,}24$~млрд~руб.
        \item DM~$= 6{,}247$, $p<0{,}001$ (от LGB-компоненты, доля 0,75)
      \end{itemize}
    \end{block}
    \vspace{0.15em}
    {\footnotesize\color{hsedark}\textbf{SHAP:} монетарные переменные \textbf{доминируют}; ДКП ЦБ РФ однородна по регионам.}
  \end{column}
  \begin{column}{0.43\textwidth}
    \textbf{\color{hsedark}Дополнительные выводы:}
    \begin{itemize}\setlength{\itemsep}{0pt}
      \small
      \item Ablation: наибольший прирост~--- производственный блок (A0$\to$A1)
      \item Robustness checks подтвердили устойчивость ранжирования
      \item GPBoost нестабилен OOS~--- GP-компонента отключается
    \end{itemize}
    \vspace{0.05em}
    \textbf{\color{hsedark}Ограничения:}
    \begin{itemize}\setlength{\itemsep}{0pt}
      \small
      \item Малая панель ($\approx\!20$~набл./регион)
      \item Часть данных доступна с лагом публикации
      \item SHAP объясняет, но \textit{не доказывает} причинность
    \end{itemize}
    \vspace{0.05em}
    {\footnotesize\textbf{\color{hsedark}Практ. значимость:} nowcast ВРП регионов до публикации Росстата.}
  \end{column}
\end{columns}
\end{frame}

% ══════════════════════════════════════════════════════════════════════════════
% СЛАЙД 17: Спасибо за внимание
% ══════════════════════════════════════════════════════════════════════════════
\begin{frame}[plain]
\begin{tikzpicture}[remember picture, overlay]
  \fill[hsedark]
    (current page.north west) rectangle
    ([yshift=-2.2cm]current page.north east);
  \fill[hsered]
    ([yshift=-2.20cm]current page.north west) rectangle
    ([yshift=-2.38cm]current page.north east);
  \fill[hsedark]
    (current page.south west) rectangle
    ([yshift=0.72cm]current page.south east);
  \draw[white, line width=1.8pt]
    ([xshift=1.1cm, yshift=-1.1cm]current page.north west) circle (0.55cm);
  \node[white, font=\bfseries\scriptsize] at
    ([xshift=1.1cm, yshift=-1.1cm]current page.north west) {ВШЭ};
  \node[white, font=\scriptsize\bfseries, anchor=west] at
    ([xshift=1.9cm, yshift=-1.1cm]current page.north west)
    {НИУ «ВЫСШАЯ ШКОЛА ЭКОНОМИКИ»};
  \node[white, font=\tiny, anchor=west] at
    ([xshift=0.5cm, yshift=0.46cm]current page.south west)
    {ФКН НИУ ВШЭ \textbar{} ОП «Анализ больших данных»};
  \node[white, font=\tiny, anchor=east] at
    ([xshift=-0.5cm, yshift=0.46cm]current page.south east)
    {Москва, 2026};
\end{tikzpicture}

\vspace{2.4cm}
\begin{center}
  {\color{hsedark}\LARGE\bfseries Спасибо за внимание!}\\[0.8em]
  {\color{hsemidgray}\rule{0.55\textwidth}{0.4pt}}\\[0.7em]
  {\small
  \begin{tabular}{r@{\hspace{0.5em}}l}
    \textbf{\color{hsedark}Автор:} &
      Семенюченко Артём Геннадьевич\\[0.2em]
    \textbf{\color{hsedark}Науч. руководитель:} &
      Паточенко Евгений Анатольевич\\[0.2em]
    \textbf{\color{hsedark}Тема:} &
      Прогнозирование ВРП субъектов РФ\\
    & ансамблевыми методами машинного обучения
  \end{tabular}}
  \\[0.8em]
  \begin{tikzpicture}
    \node[draw=hsedark, rounded corners=3pt, inner sep=5pt,
          fill=hsegray, font=\small\color{hsedark}] {
      \begin{tabular}{c}
        Лучшая модель: CB$\!\times\!$0,25 + LGB$\!\times\!$0,75\\
        \textbf{RMSE~$= 0{,}0382$} \quad \textbf{MAPE~$= 2{,}88\%$} \quad $R^2_{\text{OOS}} = 0{,}144$
      \end{tabular}};
  \end{tikzpicture}
\end{center}
\end{frame}

\end{document}
